\section{Билет 3. Метод потенциалов решения транспортных задач}

\ticketq{Метод потенциалов решения транспортных задач.}

\paragraph{Постановка.}
$m$ поставщиков с запасами $a_i$, $n$ потребителей с потребностями $b_j$,
тарифы $c_{ij}$, перевозки $x_{ij}$:
\[
  F(X) = \sum_{i=1}^{m}\sum_{j=1}^{n} c_{ij}x_{ij} \to \min , \quad
  \sum_{j=1}^{n} x_{ij} = a_i , \quad
  \sum_{i=1}^{m} x_{ij} = b_j , \quad x_{ij} \ge 0 .
\]
Задача \textbf{закрытая}, если $\sum_i a_i = \sum_j b_j$; иначе открытая и
приводится к закрытой фиктивным поставщиком ($a_{m+1} = \sum_j b_j - \sum_i a_i$)
или потребителем ($b_{n+1} = \sum_i a_i - \sum_j b_j$) с нулевыми тарифами.

\paragraph{Баланс необходим и достаточен.}
\textit{Задача разрешима $\iff$ $\sum_i a_i = \sum_j b_j$.}

\textsc{Док-во.} \emph{Необходимость:} суммируя ограничения по строкам и по
столбцам, получаем одну и ту же величину:
$\sum_i a_i = \sum_i\sum_j x_{ij} = \sum_j b_j$.
\emph{Достаточность:} при балансе план существует явно —
$x_{ij} = a_i b_j / \sum_k a_k$ удовлетворяет всем ограничениям.
$\blacksquare$

\paragraph{Базисный план и потенциалы.}
Уравнений $m+n$, одно зависимо (та же сумма по строкам и по столбцам)
$\Rightarrow$ $\rank = m+n-1$ $\Rightarrow$ невырожденный базисный план имеет
ровно $m+n-1$ заполненных клеток. Если клеток меньше, план вырожден: добавляют
базисные клетки с нулевой перевозкой, \emph{так чтобы не образовался цикл}
(иначе набор клеток перестанет быть базисом).

Строке $i$ отвечает потенциал $u_i$, столбцу $j$ — $v_j$; для заполненных
клеток $v_j - u_i = c_{ij}$. Уравнений $m+n-1$, неизвестных $m+n$
$\Rightarrow$ один потенциал задаётся произвольно ($u_1 = 0$), остальные
однозначны. Для пустых клеток $\Delta_{ij} = v_j - u_i - c_{ij}$; условие
оптимальности — $\Delta_{ij} \le 0$.

\paragraph{Теорема (об оптимальности плана).}
\textit{Если для допустимого базисного плана $X$ найдены $u_i, v_j$ с}
\[
  v_j - u_i = c_{ij} \ \text{(заполненные)} , \qquad
  v_j - u_i \le c_{ij} \ \text{(пустые)} ,
\]
\textit{то $X$ оптимален.}

\textsc{Док-во.} Пусть $Y = (y_{ij})$ — любой допустимый план. Для всех клеток
$c_{ij} \ge v_j - u_i$; умножая на $y_{ij} \ge 0$ и суммируя:
\[
  F(Y) \ \ge\ \sum_{i}\sum_{j}(v_j - u_i)y_{ij}
       = \sum_{j} v_j \underbrace{\sum_{i} y_{ij}}_{=\,b_j}
       - \sum_{i} u_i \underbrace{\sum_{j} y_{ij}}_{=\,a_i}
       = \sum_j v_j b_j - \sum_i u_i a_i .
\]
Для самого $X$: на заполненных клетках $c_{ij} = v_j - u_i$, на пустых
$x_{ij} = 0$ $\Rightarrow$ замена $c_{ij} \to v_j - u_i$ законна во всех
клетках, и
\[
  F(X) = \sum_i\sum_j (v_j - u_i)x_{ij} = \sum_j v_j b_j - \sum_i u_i a_i
  \ \Rightarrow\ F(Y) \ge F(X) . \qquad \blacksquare
\]

Величина $\sum_j v_j b_j - \sum_i u_i a_i$ — значение двойственной задачи:
неравенство есть слабая двойственность, совпадение на $X$ — сильная, а проверка
$\Delta_{ij} \le 0$ есть проверка двойственной допустимости.

\paragraph{Цикл пересчёта.}
Если $\exists\,(p,q)$ пустая с $\Delta_{pq} > 0$, она вводится в базис. Цикл
пересчёта — замкнутая ломаная: начинается в $(p,q)$, идёт по строкам и
столбцам, поворачивает только в заполненных клетках, возвращается в $(p,q)$.
Знаки чередуются от $\oplus$ в $(p,q)$: $C_+$ и $C_-$ — клетки со знаками
$\oplus$ и $\ominus$. Величина пересчёта и новый план:
\[
  \theta = \min_{(i,j) \in C_-} x_{ij} ,
  \qquad
  x_{ij} := x_{ij} + \theta \ \ \text{на } C_+ , \qquad
  x_{ij} := x_{ij} - \theta \ \ \text{на } C_- .
\]
Клетка из $C_-$, на которой достигнут минимум, выходит из базиса.

\paragraph{Теорема (об улучшении плана).}
\textit{$\Delta F = -\theta\,\Delta_{pq} < 0$.}

\textsc{Док-во.}
$\Delta F = \theta\big(\sum_{C_+} c_{ij} - \sum_{C_-} c_{ij}\big)$.
Все клетки цикла, кроме $(p,q)$, заполнены $\Rightarrow$ для них
$c_{ij} = v_j - u_i$; для стартовой $c_{pq} = v_q - u_p - \Delta_{pq}$. Цикл
замкнут: каждая задействованная строка и каждый столбец встречаются ровно
дважды — по разу со знаками $\oplus$ и $\ominus$ $\Rightarrow$ все $u_i, v_j$
взаимно сокращаются, остаётся добавка стартовой клетки:
\[
  \sum_{C_+} c_{ij} - \sum_{C_-} c_{ij} = -\Delta_{pq}
  \ \Rightarrow\ \Delta F = -\theta\,\Delta_{pq} < 0
  \quad (\theta > 0,\ \Delta_{pq} > 0) . \qquad \blacksquare
\]
Отсюда правило выбора: брать клетку с \emph{наибольшим} $\Delta_{ij} > 0$.

\paragraph{Алгоритм метода потенциалов.}
\begin{csteps}
  \item Проверить закрытость, при необходимости ввести фиктивного поставщика
        или потребителя.
  \item Построить начальный базисный план (северо-западный угол либо
        минимальный элемент); проверить число клеток $m+n-1$.
  \item Найти потенциалы: $u_1 = 0$, далее $v_j - u_i = c_{ij}$ по заполненным
        клеткам.
  \item Для пустых клеток вычислить $\Delta_{ij} = v_j - u_i - c_{ij}$.
  \item Все $\Delta_{ij} \le 0$ — план оптимален, конец.
  \item Иначе взять клетку с максимальным $\Delta_{ij} > 0$, построить цикл,
        найти $\theta$, пересчитать план, вывести клетку с нулём на минусе;
        к~шагу~3.
\end{csteps}

Метод потенциалов — симплекс-метод в табличной форме: $u_i, v_j$ играют роль
потенциалов $y^k$, $\Delta_{ij}$ — роль оценок $-d[j]$.

\seefull{гл.~2, \S4, \S5, \S7 (численный пример в \S4)}
